% !TEX program = xelatex
\documentclass[UTF8,11pt,a4paper]{ctexart}

\usepackage[a4paper,margin=2.2cm]{geometry}
\usepackage{amsmath,amssymb,amsfonts,mathtools}
\usepackage{booktabs}
\usepackage{multirow}
\usepackage{siunitx}
\usepackage{graphicx}
\usepackage{enumitem}
\usepackage{hyperref}
\usepackage{microtype}
\usepackage{tcolorbox}
\usepackage{algorithm}
\usepackage{algpseudocode}
\usepackage{caption}
\usepackage{array}

\hypersetup{colorlinks=true,linkcolor=blue,urlcolor=blue,citecolor=blue}
\setlist[itemize]{topsep=2pt,itemsep=2pt,parsep=0pt,leftmargin=1.6em}
\setlist[enumerate]{topsep=2pt,itemsep=2pt,parsep=0pt,leftmargin=1.6em}

\title{在 V3 基础上的改进方案：\,\textbf{事件分布生成策略}驱动的课程/对抗训练（V3+EDRL）}
\author{（面向：二值动作\,\&\,二值奖励的 ALNS--RL 联合系统）}
\date{\today}

\begin{document}
\maketitle

\begin{abstract}
针对现有 V3 在纯 OOD（训练仅见分布 $A$，部署切换到未见分布 $B$ 且测试期不可更新参数）场景下的性能上限与动作塌缩问题，本文给出一个在工程上可落地、与现有``不确定性事件文件''生成机制完全兼容的改进框架：\textbf{V3+EDRL}（Environment-Designer RL）。该框架在 V3 主策略（内层）之外引入一个\textbf{事件分布生成策略}（外层），外层策略输出不确定性事件生成参数（如均值组合、占比、切换模式等），生成事件文件并驱动主策略训练。外层奖励采用``难但可学''（hard-but-learnable）原则：既推动主策略暴露薄弱点，又避免退化到不可辨识噪声；同时显式加入\textbf{反塌缩约束}，确保动作 1 不被样本饥饿。我们给出严格的数学形式化（双层优化/对抗课程）、完整训练分期（Phase1--Phase4）、可复现的配置建议与伪代码，并分析该机制为何能够在不改变 V3 编码归纳偏置的前提下，主要从\textbf{数据分布覆盖与样本结构}层面缓解 OOD 失败。
\end{abstract}

\tableofcontents
\newpage

\section{问题背景与符号约定}

\subsection{任务结构（与你们系统口径一致）}
系统在每个 episode（或一个可重复的交互单元）中，主策略（记为``内层''策略）与 ALNS/环境进行交互，决策为二值动作
\begin{equation}
 a_t \in \{0,1\},
\end{equation}
并获得二值奖励
\begin{equation}
 r_t \in \{0,1\}.
\end{equation}
环境包含由不确定性事件驱动的随机因素。当前工程实现中，不确定性事件在运行前由脚本生成并写入文件（例如 CSV/JSON），主程序读取该文件驱动后续交互。

\subsection{分布键与 OOD 模式}
用 $\mathcal{D}(\phi)$ 表示事件分布族，$\phi$ 是可配置参数（例如均值、方差、切换模式、占比、种子等）。你们现有的 distribution key（如 $O\_{10,90}$、$F2\_{10,60}$、$G\_{10,90,60}$ 等）可视作对 $\phi$ 的离散化封装。

纯 OOD 的典型设定（对应你们的 $O$ 类 key）：
\begin{itemize}
\item 训练：采样 $\phi=\phi_A$（如均值 10）；
\item 测试：切换到未见的 $\phi=\phi_B$（如均值 60/90）；
\item 测试期：不允许更新主策略参数。
\end{itemize}

\subsection{V3 的输入结构化编码}
V3 的核心归纳偏置是将输入分解为
\begin{equation}
 x_t = [o_t,\, m_t,\, \Delta o_t],\quad \Delta o_t=o_t-o_{t-1},
\end{equation}
并采用分路编码/聚合（见第\ref{sec:v3}节）。该结构通常能在训练分布内抑噪、提高样本效率，但在纯 OOD 下仍可能失败，尤其在二值奖励+极不平衡样本时容易出现动作塌缩。

\section{V3 基线：结构化编码 + PPO 目标}
\label{sec:v3}

\subsection{网络结构（抽象表达）}
将三路输入分别编码：
\begin{equation}
 h^{(o)}_t=f_o(o_t),\quad h^{(m)}_t=f_m(m_t),\quad h^{(\Delta)}_t=f_\Delta(\Delta o_t),
\end{equation}
聚合得到表征
\begin{equation}
 h_t = \textsc{Agg}\big(h^{(o)}_t, h^{(m)}_t, h^{(\Delta)}_t\big),
\end{equation}
并输出策略与价值：
\begin{equation}
 \pi_\theta(a_t\mid x_t)=\textsc{Softmax}(W_\pi h_t),\qquad V_\theta(x_t)=W_V h_t.
\end{equation}

\subsection{PPO 目标（你们工程可直接复用）}
记优势估计为 $\hat A_t$，重要性比率为 $r_t(\theta)=\frac{\pi_\theta(a_t\mid x_t)}{\pi_{\theta_{\mathrm{old}}}(a_t\mid x_t)}$。PPO 的 clipped surrogate：
\begin{equation}
\label{eq:ppo}
\mathcal{L}^{\mathrm{CLIP}}(\theta)=\mathbb{E}_t\Big[\min\big(r_t(\theta)\hat A_t,\, \mathrm{clip}(r_t(\theta),1-\epsilon,1+\epsilon)\hat A_t\big)\Big].
\end{equation}
总体损失：
\begin{equation}
\mathcal{L}(\theta)= -\mathcal{L}^{\mathrm{CLIP}}(\theta) + c_v\,\mathcal{L}^{V}(\theta) - c_e\,\mathcal{H}(\pi_\theta),
\end{equation}
其中 $\mathcal{H}$ 为熵项。

\section{为何纯 OOD 下仍失败：动作塌缩与信息-样本结构瓶颈}

\subsection{动作塌缩的自增强机制}
在二值动作/二值奖励设定中，若训练分布 $\phi_A$ 下 action=0 的即时回报显著更高（或方差更低），策略会快速把 $\pi(a=1\mid x)$ 压到极小。之后采样中 action=1 变得稀缺，导致：
\begin{itemize}
\item action=1 的正样本更稀缺；
\item 价值/优势估计对 action=1 更不稳定；
\item 梯度进一步推动 action=1 概率下降。
\end{itemize}
这一反馈环会把策略锁死在 action=0，形成你们观察到的 ``action1\_rate $\approx 0$'' 现象。

\subsection{``补数据''何时有效、何时无效}
直观上，若仅增加训练量但仍只采样 $\phi_A$，主策略只会更确信``做 0''；因此需要让训练阶段覆盖 $\phi_B$ 或其邻域。

但要注意：\textbf{覆盖并不等于可学}。若在 $\phi_B$ 下，当前可观测输入 $x_t$ 对奖励差异几乎不可辨识（例如 $P(r=1\mid a=1,x)\approx P(r=1\mid a=0,x)$ 在大部分 $x$ 上成立），则再多数据也只能学到常数策略或随机化。V3+EDRL 的设计原则之一，就是让外层策略倾向生成 ``难但可学'' 的分布，而不是纯噪声。

\section{V3+EDRL：外层事件分布生成策略驱动的课程/对抗训练}

\subsection{核心思想}
保持 V3 主策略架构不变，通过外层策略 $\pi_\psi$ 输出事件生成参数 $a_g$，生成事件文件 $E$，驱动内层训练。

\begin{tcolorbox}[title=一句话定义,width=\textwidth,colback=gray!4]
\textbf{V3+EDRL}：\;在不改变 V3 编码归纳偏置的前提下，引入外层 RL 在``事件分布参数空间''上做自动课程/受约束对抗，使内层主策略在训练期就看到目标 OOD 分布及其关键样本结构，从而缓解动作塌缩与分布覆盖不足。
\end{tcolorbox}

\subsection{事件分布参数化与文件生成}
用 $a_g$ 表示外层动作（即事件生成参数）。推荐一个最小但足够表达的动作空间：
\begin{equation}
\label{eq:ag}
 a_g=(\mu_1,\mu_2,p,\text{pattern},\sigma,\text{seed}).
\end{equation}
其中：
\begin{itemize}
\item $\mu_1,\mu_2\in\{10,30,60,90\}$（或你们已有候选集）；
\item $p\in[0,1]$ 表示前 $p$ 比例使用均值 $\mu_1$、后 $1-p$ 使用 $\mu_2$；
\item $\text{pattern}\in\{ab,aba,abba,\text{random\_mix}\}$ 与你们 key 体系对齐；
\item $\sigma$ 控制波动强度（可先固定）；
\item seed 仅为可复现。
\end{itemize}

定义事件文件生成器 $G$：
\begin{equation}
 E=G(a_g),
\end{equation}
其中 $E$ 是磁盘文件（或其内存表示），包含每个时间步/任务点的不确定性事件参数。

\subsection{内层与外层目标：双层优化视角}
内层主策略参数为 $\theta$，外层生成策略参数为 $\psi$。

\paragraph{内层目标（主策略）}
在给定事件文件 $E$ 下，主策略的回报为
\begin{equation}
 J(\theta;E)=\mathbb{E}\Big[\sum_{t=0}^{T-1}\gamma^t r_t\;\big|\;E,\pi_\theta\Big].
\end{equation}
内层训练用 PPO（式\eqref{eq:ppo}）更新 $\theta$。

\paragraph{外层目标（生成策略）}
外层的目标不是单纯``让内层更差''（那会生成不可学噪声），也不是单纯``让内层更好''（那会回到最易分布）。推荐采用\textbf{难但可学（hard-but-learnable）}奖励：
\begin{equation}
\label{eq:rg}
 r_g(a_g)=\underbrace{-J(\bar\theta;G(a_g))}_{\text{难度}}+\lambda\underbrace{\Delta J(\bar\theta;G(a_g))}_{\text{可学习性}}+\eta\underbrace{B(G(a_g))}_{\text{反塌缩}}-\gamma\underbrace{D(a_g)}_{\text{可行域约束}}.
\end{equation}
其中 $\bar\theta$ 是冻结的主策略（Phase2），\;$\Delta J$ 定义为在同一 $E$ 上进行少量内层更新后的提升幅度：
\begin{equation}
\Delta J(\bar\theta;E)=J(\theta^+;E)-J(\bar\theta;E),
\end{equation}
$\theta^+$ 是从 $\bar\theta$ 复制出来、在 $E$ 上做 $K$ 次 PPO 更新得到的临时参数。

\paragraph{反塌缩项 $B(E)$}
令在事件文件 $E$ 下采样得到的 action1 比例为 $\hat\rho_1(E)$。可用一个温和的 ``贴近目标采样率'' 形式：
\begin{equation}
 B(E)=-\big|\hat\rho_1(E)-\rho^*\big|,
\end{equation}
其中 $\rho^*$ 推荐设为 $5\%\sim 15\%$（视任务而定）。也可替换为策略熵 $\mathcal{H}$ 的下界惩罚。

\paragraph{约束项 $D(a_g)$}
用于限制外层生成的参数不越界，防止奖励黑客（例如极端参数导致文件异常）。最简单做法是把硬约束转为罚项：
\begin{equation}
 D(a_g)=\mathbb{I}[a_g\notin\mathcal{A}_{\mathrm{valid}}]\cdot M + \text{soft\_penalty}(a_g),
\end{equation}
$M$ 为大常数。

\subsection{外层算法选择}
外层动作空间可离散（枚举 key）或连续（$p$ 连续）。两种推荐：
\begin{itemize}
\item \textbf{离散版（最稳，优先建议）}：外层直接输出 distribution key（如从 $\{O\_*,G\_*,F\_*,R\_*\}$ 选一个），这等价于在 key 空间做 bandit/PPO。
\item \textbf{参数版（更灵活）}：外层输出 $(\mu_1,\mu_2,p,\text{pattern})$，其中 $p$ 连续可用 Beta 分布参数化或直接输出 $p\in(0,1)$。
\end{itemize}

\section{训练分期（Phase1--Phase4）}

\subsection{Phase1：V3 主策略基线训练}
按现有流程，在训练分布（如均值 10 或你们既定的训练 key 集）上训练 V3，得到基线参数 $\theta_0$，并保存所有可复现配置（包括事件文件种子、分布 key、模型结构、PPO 超参）。

\subsection{Phase2：训练外层生成策略（冻结内层）}
冻结内层为 $\bar\theta\leftarrow\theta_0$。每次外层提出 $a_g$，生成事件文件 $E$，评估 $J(\bar\theta;E)$，再用临时复制参数 $\theta^+$ 估计 $\Delta J$，从而得到外层奖励 $r_g$（式\eqref{eq:rg}）。外层用 PPO/REINFORCE 更新 $\psi$。

\subsection{Phase3：联合课程训练（外层生成 + 内层可更新）}
用外层策略生成事件文件驱动内层训练，但为避免内层被带偏，建议混合数据源：
\begin{equation}
 E \sim \alpha\,G(\pi_\psi) + (1-\alpha)\,\mathcal{D}_{\mathrm{base}},
\end{equation}
其中 $\alpha$ 采用课程式递增（如 $0.2\to 0.5\to 0.8$）。同时维护一个``难例池''（外层历史上高 $r_g$ 的若干 $a_g$）进行 replay。

\subsection{Phase4：测试（与现有口径一致）}
冻结内层参数，按你们既定的 O/G/F key 集进行评估：
\begin{itemize}
\item OOD：$O\_{10,60}$、$O\_{10,90}$ 等；
\item 泛化：$G\_{10,90,60}$ 等；
\item 遗忘：$F1/F2$；
\item 关键诊断指标：avg\_reward、action1\_rate、p\_action1、奖励方差、熵等。
\end{itemize}

\section{算法伪代码（可直接交给 Codex 实施）}

\begin{algorithm}[H]
\caption{V3+EDRL（概念版伪代码）}
\begin{algorithmic}[1]
\State \textbf{Input:} 基础训练分布集合 $\mathcal{D}_{\mathrm{base}}$，事件生成器 $G(\cdot)$，测试 key 集 $\mathcal{K}_{\mathrm{test}}$
\State \textbf{Phase1:} 训练内层 V3（PPO）得到 $\theta_0$
\State \textbf{Phase2:} 冻结 $\bar\theta\leftarrow\theta_0$；初始化外层参数 $\psi$
\For{$i=1$ to $N_{\mathrm{outer}}$}
  \State 采样外层动作 $a_g\sim\pi_\psi(\cdot\mid s_g)$
  \State 生成事件文件 $E\leftarrow G(a_g)$；若非法则给惩罚并 continue
  \State 用冻结内层评估 $J\leftarrow J(\bar\theta;E)$
  \State 复制临时参数 $\theta\_\mathrm{tmp}\leftarrow\bar\theta$；在 $E$ 上做 $K$ 次 PPO 更新得到 $\theta^+$
  \State 评估 $J^+\leftarrow J(\theta^+;E)$；计算 $\Delta J\leftarrow J^+-J$
  \State 计算 $\hat\rho\_1(E)$；得到外层奖励 $r_g$（式\eqref{eq:rg}）
  \State 用 $r_g$ 更新外层策略参数 $\psi$
\EndFor
\State \textbf{Phase3:} 联合训练：每轮按混合比例 $\alpha$ 从 $G(\pi_\psi)$ 与 $\mathcal{D}_{\mathrm{base}}$ 取事件文件，更新内层 $\theta$
\State \textbf{Phase4:} 冻结内层，在 $\mathcal{K}_{\mathrm{test}}$ 上评估并输出对比报告
\end{algorithmic}
\end{algorithm}

\section{参数设置建议（可复现口径）}

\subsection{内层 V3（PPO）推荐超参}
下表给出在你们二值动作/奖励、易塌缩场景下更稳的默认值（起点）。具体仍需按 R、batch size、episode 长度微调。

\begin{table}[H]
\centering
\caption{内层 PPO 建议超参（起点）}
\begin{tabular}{llp{7.5cm}}
\toprule
参数 & 建议值 & 说明\\
\midrule
学习率 $\eta$ & $3\times 10^{-4}$ & 若塌缩严重，可降至 $1\times 10^{-4}$；或用线性退火。\\
clip $\epsilon$ & 0.2 & 二值动作场景通常足够。\\
GAE $\lambda$ & 0.95 & 标准设置。\\
折扣 $\gamma$ & 0.99 & 视 episode 长度调整。\\
value coef $c_v$ & 0.5 & 标准设置。\\
entropy coef $c_e$ & 0.01（可退火） & \textbf{反塌缩关键}：前期可提高到 0.02--0.05，后期退火到 0.005。\\
min action1 rate & 5\%（训练期） & 工程上可用采样/约束实现，避免 action1 饿死（只在训练期用）。\\
\bottomrule
\end{tabular}
\end{table}

\subsection{外层生成策略（EDRL）推荐超参}
\begin{table}[H]
\centering
\caption{外层 EDRL 建议超参（起点）}
\begin{tabular}{llp{7.5cm}}
\toprule
参数 & 建议值 & 说明\\
\midrule
外层算法 & PPO（离散 key 优先） & 离散 key 版本最稳；若用连续 $p$，可用 Beta 参数化。\\
外层 rollout & 1 个 $a_g$ 对应 1 个短评估 & 外层 step 很贵，建议用短评估代理。\\
$K$（可学习性步数） & 1--3 次内层更新 & 只为估计 $\Delta J$ 的方向性，避免 Phase2 过慢。\\
$\lambda$ & 1.0 & 可学习性权重；若外层偏向纯对抗可增大。\\
$\eta$ & 0.2--0.5 & 反塌缩权重；过大可能迫使外层生成``硬塞 action1''的分布。\\
$\rho^*$ & 5\%--15\% & action1 目标采样率（训练期）。\\
$\gamma$（罚项） & 10--100 & 非法参数的惩罚强度。\\
动作空间 & $(\mu_1,\mu_2,p,\text{pattern})$ & 先固定 $\sigma$，减少自由度。\\
\bottomrule
\end{tabular}
\end{table}

\subsection{Phase3 混合比例（课程）}
建议 $\alpha$ 三阶段递增：
\begin{equation}
\alpha \in \{0.2,0.5,0.8\},
\end{equation}
每个阶段持续若干 epoch（例如按你们现有训练轮数均分）。若观察到内层在 base 分布上明显退化，则降低 $\alpha$ 或增大 replay（难例池 + base 池）。

\section{机理分析：为什么 V3+EDRL 有望打破 OOD 上限}

\subsection{从``模型改进''转为``训练分布设计''}
V3 已经给了强归纳偏置；若 OOD 失败主要来自训练分布覆盖不足与样本结构偏置（action1 饿死），继续堆网络改动往往收益有限。EDRL 把优化重点移到 ``训练数据是从哪里来、难度如何自适应升级''。

\subsection{难但可学：避免对抗退化到不可辨识噪声}
纯对抗会选择最坏分布，但最坏分布可能是不可学习的（近似随机）。引入 $\Delta J$ 明确鼓励 ``在该分布上做少量训练能变好''，使外层倾向生成具有训练价值的挑战分布。

\subsection{反塌缩：打断动作自锁死循环}
显式约束 action1 采样率或熵下界，使内层在训练阶段持续看到 action1 的正/负样本，从而让优势估计对 action1 不至于完全失真，梯度不会永久把 action1 压到 0。

\subsection{课程与 replay：稳定双层非平稳性}
外层与内层同时学习本质是非平稳博弈。通过 Phase2 冻结内层先学习 ``会打中弱点'' 的分布，再在 Phase3 用混合比例与 replay 逐步引入，可显著提高稳定性与可复现性。

\section{工程落地要点：事件文件系统兼容与可复现}

\subsection{接口建议（不破坏现有流程）}
保留 ``生成事件文件 $\to$ 读取文件 $\to$ 训练/交互'' 的主流程，只新增：
\begin{itemize}
\item 一个外层策略输出 $a_g$；
\item 一个包装器把 $a_g$ 写成事件生成器可读的配置（YAML/JSON）；
\item 事件文件缓存（同一 $a_g$ 重复命中则复用文件路径）；
\item meta.json 里记录 $a_g$、hash、生成器版本号。
\end{itemize}

\subsection{合法性校验（必须做）}
生成事件文件后至少校验：
\begin{itemize}
\item 事件数量、时间长度与主程序预期一致；
\item 均值/方差/占比符合 $\mathcal{A}_{\mathrm{valid}}$；
\item 切换点存在且在范围内；
\item 文件解析无异常；
\item 统计摘要（均值、分位数、段内均值）与 $a_g$ 一致。
\end{itemize}
非法直接大惩罚并丢弃，避免外层通过异常文件 ``黑客'' 奖励。

\section{实现说明（结合当前项目代码结构）}

\subsection{代码级插入点（已与现有工程对齐）}
基于当前仓库，V3+EDRL 不需要推翻主干训练链路，而是复用已有入口与日志体系。关键事实如下：
\begin{itemize}
\item 统一入口在 \texttt{codes/Dynamic\_master34959.py}：\texttt{run\_generator()} 负责生成事件文件，\texttt{run\_single()} 负责设置 \texttt{RL\_ALGO\_VERSION}、\texttt{RL\_STAGE\_MODE}、\texttt{DYNAMIC\_DATA\_ROOT} 并启动 RL+ALNS。
\item 事件生成器在 \texttt{codes/generation/generate\_mixed\_parallel.py}，并且当前有 \texttt{EXPECTED\_TOTAL\_FILES=1000} 的硬覆盖逻辑，这与 ``外层小批量（3--50）'' 目标冲突。
\item 读取链路在 \texttt{codes/core/dynamic\_RL34959.py} 与 \texttt{codes/core/Intermodal\_ALNS34959.py} 的 \texttt{resolve\_dynamic\_data\_path()}，两者都优先读取环境变量 \texttt{DYNAMIC\_DATA\_ROOT}。
\item 元信息注入点在 \texttt{dynamic\_RL34959.py:get\_state()}：每次决策读取 Excel 的 \texttt{\_\_meta\_\_} sheet，提取 \texttt{gt\_mean} 与 \texttt{phase\_label}，并写入 trace/decision 日志。
\item V3 表征提取器为 \texttt{algorithms/ppo\_new/nn/context\_extractor.py:StackedContextExtractor}，输出维度 \texttt{out\_dim=64}；当前 \texttt{dynamic\_RL34959.py} 已通过 \texttt{\_decision\_extract\_h64()} 进行导出。
\item 决策级日志已存在：\texttt{rl\_decision.csv} + \texttt{h\_dump.npz}，字段由 \texttt{DECISION\_FIELDS} 定义，主键是 \texttt{decision\_id = run\_id|phase|decision\_seq}。
\end{itemize}

\subsection{为什么建议 ``新建生成器'' 而非直接改旧生成器}
旧生成器 \texttt{generate\_mixed\_parallel.py} 的设计目标是一次性物理隔离生成 1000 个文件，并服务原始表号区间（训练 0--799，测试 999--800 的倒序读表）。外层 EDRL 需要的是：
\begin{equation}
a_g=(\mu_A,\mu_B,p,n,s),\quad \mu_A,\mu_B\in\{10,30,60,90\},\;p\in(0,1),\;n\in[3,50],
\end{equation}
其中 $n$ 很小且每轮可变。为避免破坏旧实验可复现性，建议保留旧生成器不动，新增一个专用小批量生成器。

\subsection{建议新增模块与职责}
\begin{table}[H]
\centering
\caption{V3+EDRL 的最小新增模块（建议）}
\begin{tabular}{p{4.2cm}p{3.2cm}p{7.4cm}}
\toprule
模块路径（建议） & 类型 & 职责 \\
\midrule
\texttt{codes/generation/generate\_outer\_small\_batch.py} & 新增脚本 & 接收外层动作参数 $(\mu_A,\mu_B,p,n,s)$，生成每轮小批量事件文件；写入 \texttt{\_\_meta\_\_}（\texttt{gt\_mean}、\texttt{phase\_label}）与清单文件。\\
\texttt{codes/generation/outer\_batch\_schema.py} & 新增模块（可选） & 统一参数校验、分段规则、随机种子策略，避免训练期和复现实验口径漂移。\\
\texttt{codes/outer\_rl/run\_edrl\_phase2.py} & 新增脚本（调度层） & 管理 ``外层动作 $\to$ 生成数据 $\to$ 内层 train\_only 训练'' 闭环，落盘每轮 manifest 与指标。\\
\texttt{codes/analysis/validate\_outer\_batch.py} & 新增脚本 & 校验小批量文件完整性、分布占比误差、meta 一致性。\\
\bottomrule
\end{tabular}
\end{table}

\subsection{小批量生成器的核心实现点（必须落地）}
\paragraph{1) 动作到文件段的映射}
外层动作为 $(\mu_A,\mu_B,p,n,s)$，建议映射为
\begin{equation}
n_A=\mathrm{clip}\!\left(\left\lfloor p n + 0.5\right\rfloor,\,1,\,n-1\right),\qquad n_B=n-n_A,
\end{equation}
并按 ``前 $n_A$ 个文件为 A、后 $n_B$ 个文件为 B'' 写入，保证与用户指定的 ``前后占比'' 一致。

\paragraph{2) 输出目录结构（每轮独立）}
建议落盘到
\begin{equation}
\texttt{run\_dir/post\_stage/outer\_batches/iter\_k/Rr/Intermodal\_EGS\_data\_dynamic\_congestion\{i\}.xlsx},
\end{equation}
并附带 \texttt{manifest.json}（记录 $a_g$、seed、文件数、均值统计、生成时间、hash）。

\paragraph{3) 与现有读取器的契约保持不变}
每个 Excel 仍保持原 sheet 命名与字段；特别是必须写入 \texttt{\_\_meta\_\_}：
\begin{itemize}
\item \texttt{Property=gt\_mean, Value=\{mu\_A or mu\_B\}}
\item \texttt{Property=phase\_label, Value=\{A or B\}}
\end{itemize}
这样 \texttt{dynamic\_RL34959.py:get\_state()} 可零改动读取并写入 RL trace/decision 统计。

\subsection{小批量文件与 ``大表号'' 的兼容映射（关键难点）}
当前 ALNS 主循环使用固定表号制度（训练上限 799，测试区间 999$\rightarrow$800）。若 Phase2 只生成 $n\ll1000$ 文件，读取层必须加入表号映射：
\begin{equation}
i(t)=\big(t-t_0\big)\bmod n,\qquad i(t)\in\{0,\dots,n-1\},
\end{equation}
其中 $t$ 为当前 \texttt{table\_number}，$t_0$ 为该阶段起点。实现上建议仅在 \texttt{resolve\_dynamic\_data\_path()} 增加 ``映射模式''，不改 ALNS 其余逻辑。

这一步是 V3+EDRL 能否工程可用的分水岭：若不加映射，系统会尝试访问不存在的 \texttt{congestion800+} 文件并触发空读/异常。

\subsection{Phase 调度与协议约束（不改 PPO 主损失）}
结合现有 \texttt{stage\_mode} 机制（\texttt{train\_eval/train\_only/eval\_only}），建议执行顺序：
\begin{enumerate}
\item Phase1：\texttt{train\_only} 在训练分布完成 V3 基线训练，保存 \texttt{policy\_theta0}。
\item Phase2：仍用 \texttt{train\_only}，但把 \texttt{DYNAMIC\_DATA\_ROOT} 指向外层生成的小批量目录，进行后训练/微调；此阶段不访问最终测试分布 B 的 reward。
\item Phase3：\texttt{train\_only} 或 \texttt{train\_eval}，执行 ``外层生成 + 内层可更新'' 的联合课程训练（混合 base/outer 数据并可启用 replay）。
\item Phase4：\texttt{eval\_only} 在最终测试分布执行只预测评测，不做参数更新。
\end{enumerate}
该方案保持 ``测试阶段不更新参数'' 约束：仅 Phase4 为纯评测，Phase1--Phase3 均允许训练更新。

\subsection{可直接复用的日志与验收链路}
当前项目已具备决策级与表征导出能力，可直接作为 Phase2 数据质量控制：
\begin{itemize}
\item \texttt{rl\_decision.csv}：每条决策含 \texttt{decision\_id}、\texttt{phase}、\texttt{action}、\texttt{reward}、\texttt{phase\_label}、\texttt{h\_index}。
\item \texttt{h\_dump.npz}：\texttt{decision\_id[N]} 与 \texttt{h[N,64]}。
\item \texttt{codes/analysis/validate\_decision\_artifacts.py}：已有 \texttt{pairing\_rate}、\texttt{h\_align\_rate}、按组统计 (\texttt{action1\_rate}, $E[r|a]$) 的验收框架。
\end{itemize}
建议在论文与工程验收中固定阈值（例如 pairing/h\_align 的下限）并将失败 run 标记为无效样本，避免把数据对齐问题误判为算法问题。

\subsection{实施优先级（建议顺序）}
\begin{enumerate}
\item 先实现并验证 ``小批量生成器 + 表号映射 + manifest''（纯数据层改动）。
\item 再接入外层策略（先离散 key，后连续参数 $p$）。
\item 最后开展完整 V3+EDRL 双层训练实验与消融。
\end{enumerate}
该顺序可最大化复用当前代码并最小化一次性改动风险。

\section{日志与可复现规范执行计划（外层 EDRL 重点）}

\subsection{设计目标与约束}
围绕你提出的核心风险（``外层文件生成了，但主程序读错路径''），日志与目录规范必须满足以下约束：
\begin{itemize}
\item \textbf{单轮可回放}：任意一轮 outer iter 都能通过 manifest 复原 ``动作参数 $\to$ 文件集合 $\to$ 内层训练输入''。
\item \textbf{路径强一致}：生成路径、运行时读取路径、日志记录路径三者必须同源并可交叉校验。
\item \textbf{失败可定位}：任何读文件失败都能定位到 ``哪一轮、哪一个表号、映射到哪一个文件、为何失败''。
\item \textbf{兼容现有主干}：不破坏现有 \texttt{run\_single()}、\texttt{DYNAMIC\_DATA\_ROOT}、\texttt{rl\_trace.csv}、\texttt{rl\_decision.csv} 体系。
\end{itemize}

\subsection{建议目录规范（每个 run 内）}
现有系统已把 run 根目录统一到 \texttt{codes/logs/<run\_id>/}。建议在其下新增固定结构：
\begin{equation}
\texttt{post\_stage/outer\_batches/iter\_k/Rr/...xlsx}
\end{equation}
并配套如下文件：
\begin{table}[H]
\centering
\caption{外层阶段建议产物（最小闭环）}
\begin{tabular}{p{5.4cm}p{10.0cm}}
\toprule
文件路径 & 用途 \\
\midrule
\texttt{post\_stage/manifest.json} & 记录总配置：base run、算法版本、阶段开关、路径模式、索引模式、随机种子。 \\
\texttt{post\_stage/outer\_actions.csv} & 每轮一行：\texttt{iter\_id, mu\_a, mu\_b, p, n\_files, seed, pattern, objective\_score}。 \\
\texttt{post\_stage/outer\_generation.csv} & 每轮每文件一行：\texttt{iter\_id, logical\_idx, file\_path, gt\_mean, phase\_label, hash, ok}。 \\
\texttt{post\_stage/outer\_train\_round.csv} & 每轮训练摘要：\texttt{iter\_id, ckpt\_in, ckpt\_out, learn\_steps, avg\_reward, action1\_rate, elapsed\_s}。 \\
\texttt{post\_stage/outer\_path\_map.csv} & 路径映射审计：\texttt{iter\_id, table\_number, mapped\_idx, mapped\_file, exists, read\_ok}。 \\
\texttt{post\_stage/outer\_console.log} & 控制台详细文本输出（建议和屏幕一致）。 \\
\bottomrule
\end{tabular}
\end{table}

\subsection{读取路径改造计划（避免主程序读取错误）}
当前读取函数在两个模块各有一份（\texttt{dynamic\_RL34959.py} 和 \texttt{Intermodal\_ALNS34959.py}）。执行计划是：
\begin{enumerate}
\item 统一抽象一个 ``表号到文件号'' 映射规则（直接模式 / 取模模式）。
\item 两个 \texttt{resolve\_dynamic\_data\_path()} 同步使用该规则，避免 RL 与 ALNS 视图不一致。
\item 增加严格模式：若映射后文件不存在，立刻报错并输出映射上下文，不允许静默 fallback。
\end{enumerate}

建议配置变量（计划）：
\begin{itemize}
\item \texttt{RL\_DYNAMIC\_INDEX\_MODE=direct|mod}
\item \texttt{RL\_DYNAMIC\_FILE\_COUNT=<n>}
\item \texttt{RL\_DYNAMIC\_TABLE\_BASE=<t0>}
\item \texttt{RL\_DYNAMIC\_STRICT\_PATH=1}
\item \texttt{RL\_DYNAMIC\_MANIFEST=<abs\_path>}
\end{itemize}

映射规则建议固定为：
\begin{equation}
i(t)=
\begin{cases}
t, & \text{direct}\\
\big(t-t_0\big)\bmod n, & \text{mod}
\end{cases}
\end{equation}
并将 $(t,i(t),\text{file\_path})$ 写入 \texttt{outer\_path\_map.csv}，用于复核。

\subsection{控制台超级详细输出规范（运行时一眼看懂）}
建议全程使用统一前缀标签，且每一步均打印关键参数、路径与统计。推荐标签：
\texttt{[OUTER][INIT]}, \texttt{[OUTER][ITER]}, \texttt{[OUTER][GEN]}, \texttt{[OUTER][PATH]},
\texttt{[OUTER][TRAIN]}, \texttt{[OUTER][CHECK]}, \texttt{[OUTER][DONE]}, \texttt{[OUTER][ERROR]}。

建议输出模板（示例）：
\begin{verbatim}
[OUTER][INIT] run_id=v3_edrl_O10_90_s42 base_ckpt=... stage_mode=train_only
[OUTER][INIT] data_root=.../post_stage/outer_batches index_mode=mod strict=1
[OUTER][INIT] action_space=mu in {10,30,60,90}, p in (0,1), n in [3,50]

[OUTER][ITER] iter=03/20 action=(muA=10,muB=90,p=0.30,n=12,seed=42003)
[OUTER][GEN]  iter=03 start target_dir=.../iter_003/R30
[OUTER][GEN]  iter=03 done files=12 ok=12 fail=0 hash_mismatch=0 elapsed=2.41s
[OUTER][GEN]  iter=03 stats gt_mean_A=10.3 gt_mean_B=89.7 ratio_A=0.33 ratio_B=0.67

[OUTER][PATH] iter=03 t0=400 n=12 sample_map: t=400->0 t=401->1 ... t=411->11 t=412->0
[OUTER][PATH] iter=03 strict_check missing=0 unreadable=0

[OUTER][TRAIN] iter=03 ckpt_in=...theta_iter002.zip ckpt_out=...theta_iter003.zip
[OUTER][TRAIN] iter=03 learn_steps=200 avg_reward=0.512 action1_rate=0.074 elapsed=18.6s
[OUTER][TRAIN] iter=03 rl_decision_rows=446 pairing=0.993 h_align=0.993

[OUTER][CHECK] iter=03 objective=-J+lambda*dJ+eta*B-gamma*D = 0.183 (J=0.488,dJ=0.022,B=0.061,D=0)
[OUTER][DONE]  iter=03 status=PASS cumulative_best=0.201 best_iter=2
\end{verbatim}

异常模板（必须包含定位字段）：
\begin{verbatim}
[OUTER][ERROR] iter=03 stage=PATH_CHECK table=417 mapped_idx=5
[OUTER][ERROR] expected_file=.../iter_003/R30/Intermodal_EGS_data_dynamic_congestion5.xlsx
[OUTER][ERROR] reason=FileNotFound strict=1 action=ABORT_ITER
\end{verbatim}

\subsection{截至 2026-03-01 的实施进展（已完成）}
\begin{table}[H]
\centering
\caption{V3+EDRL 当前已完成事项（代码已落地）}
\begin{tabular}{p{1.4cm}p{2.2cm}p{7.2cm}p{3.8cm}}
\toprule
阶段 & 状态 & 已完成产出 & 主要证据路径 \\
\midrule
S0 & 已完成 & 协议冻结 v1.1：统一 outer--inner 交互参数、日志 schema、内层停止模式（\texttt{fixed\_n/converge}）与目标函数口径。 & \texttt{codes/outer\_rl/run\_edrl\_phase2.py} \\
S1 & 已完成 & 小批量事件生成闭环：按 outer action 生成批次文件、写 \texttt{manifest} / \texttt{outer\_generation.csv} / hash，并支持生成后校验。 & \texttt{codes/generation/generate\_outer\_small\_batch.py}, \texttt{codes/analysis/validate\_outer\_batch.py} \\
S2 & 已完成 & 路径映射与严格读取闭环：\texttt{direct/mod} 映射、严格模式、路径审计日志 \texttt{outer\_path\_map.csv}、读通率校验。 & \texttt{codes/analysis/validate\_outer\_path\_map.py} \\
S3 & 已完成（最小可用版） & 外层调度主链路打通：\texttt{outer action $\rightarrow$ generate $\rightarrow$ inner train\_only $\rightarrow$ objective/log}，支持 \texttt{random/ucb/pg} 三种策略模式。 & \texttt{codes/outer\_rl/run\_edrl\_phase2.py} \\
S3+ & 已完成（稳定性修复） & 修复 checkpoint 退出路径丢失问题（\texttt{SystemExit(78)} 场景仍可落盘）；修正 \texttt{processed\_tables} 统计口径，降低 \texttt{fixed\_budget\_n\_partial} 误报。 & \texttt{codes/core/dynamic\_RL34959.py}, \texttt{codes/outer\_rl/run\_edrl\_phase2.py} \\
\bottomrule
\end{tabular}
\end{table}

\paragraph{当前非阻塞已知项}
退出阶段仍可能出现 \texttt{tkinter atexit} 噪声（\texttt{main thread is not in main loop}）；不影响本轮训练完成、日志落盘与 checkpoint 持久化。

\subsection{分阶段总计划（S0--S6，改进版）}
\paragraph{S0：规范冻结（2026-03-01 $\sim$ 2026-03-02，已完成）}
\textbf{目标}：统一协议，消除实现歧义。\\
\textbf{产出}：Phase 定义、\texttt{DYNAMIC\_DATA\_ROOT} 契约、状态/奖励口径、无泄漏数据分割。\\
\textbf{验收}：冻结后仅允许增量字段，不改核心语义。

\paragraph{S1：数据生成层闭环（2026-03-03 $\sim$ 2026-03-05，已提前完成）}
\textbf{目标}：小批量事件文件生成可复现、可核查。\\
\textbf{产出}：生成器、manifest、\texttt{outer\_generation.csv}、hash、一致性校验脚本。\\
\textbf{验收}：文件完整率 $100\%$；占比与均值误差在阈值内；失败可定位到 \texttt{iter/file}。

\paragraph{S2：路径映射与严格读取（2026-03-06 $\sim$ 2026-03-08，已提前完成）}
\textbf{目标}：保证主程序 ``读到正确文件''。\\
\textbf{产出}：RL/ALNS 双模块统一映射逻辑（\texttt{direct/mod}）、严格模式、\texttt{outer\_path\_map.csv} 审计。\\
\textbf{验收}：路径读通率目标 $100\%$（最低 $99\%$）；禁止静默 fallback；读错即带上下文报错。

\paragraph{S3：外层最小可用闭环（2026-03-09 $\sim$ 2026-03-12，已提前完成）}
\textbf{目标}：低风险上线外层策略与训练调度。\\
\textbf{产出}：\texttt{run\_edrl\_phase2.py} 闭环调度；\texttt{random/ucb/pg} 策略；固定预算 \texttt{fixed\_n}；日志五件套；跨 iter checkpoint 继承。\\
\textbf{验收}：多轮稳定运行；outer action 非常量退化；checkpoint 按轮可追溯。

\paragraph{S4：联合课程训练（2026-03-13 $\sim$ 2026-03-17，进行中）}
\textbf{目标}：引入课程混合与反塌缩约束，验证主收益。\\
\textbf{开发项}：实现 $\alpha$ 课程调度（建议 $0.2\rightarrow0.5\rightarrow0.8$）、难例池 replay、反塌缩监控（\texttt{action1\_rate}, $E[r|a]$, entropy）。\\
\textbf{验收}：\texttt{action1\_rate} 不塌缩；OOD 指标相对 S3 提升；base 分布退化受控（建议不超过 $5\%$）。

\paragraph{S5：对比实验与消融（2026-03-18 $\sim$ 2026-03-22，待执行）}
\textbf{目标}：形成可发表证据链。\\
\textbf{实验矩阵}：V3 baseline、V3+EDRL、去 $\Delta J$、去 $B(E)$、去 base 混合、随机 outer、固定课程基线；每组多 seed（建议 $\ge 3$）。\\
\textbf{验收}：矩阵完整可复现；统计汇总与机理解释一致（主结论在多个 seed 稳定）。

\paragraph{S6：收口与发布（2026-03-23 $\sim$ 2026-03-24，待执行）}
\textbf{目标}：工程可复现、文档可交付。\\
\textbf{产出}：一键复现实验脚本、最终配置快照、验收报告、论文图表输入表。\\
\textbf{验收}：新环境按文档可复跑；关键结果复现实验通过。

\subsection{关键里程碑（Go/No-Go 门禁）}
\begin{enumerate}
\item S1 未通过：不进入 S2。
\item S2 未通过：禁止进入任何外层训练（S3/S4 全部阻断）。
\item S3 未通过：不进入联合课程训练（S4）。
\item S4 未同时满足 ``收益 + 稳定''：回滚到 S3 调参，不进入 S5 大规模消融。
\end{enumerate}

\subsection{当前外层 RL 设计（已落地实现）}
\paragraph{目标函数与动作空间}
当前外层调度脚本 \texttt{codes/outer\_rl/run\_edrl\_phase2.py} 已实现如下目标函数（与日志字段一一对应）：
\begin{equation}
\texttt{objective}=-J+\lambda\_{\Delta J}\cdot dJ+\eta\_{\mathrm{collapse}}\cdot B-\gamma\_{\mathrm{path}}\cdot D,
\end{equation}
其中 $J\equiv\texttt{avg\_reward}$，$dJ$ 为相邻迭代的增量，$B=-|\texttt{action1\_rate}-\rho^*|$，$D$ 为路径缺失惩罚比例。

外层动作已采用参数化离散组合：
\begin{equation}
a_g=(\mu_A,\mu_B,p,n,\text{pattern},\text{seed}),
\end{equation}
其中 $(\mu_A,\mu_B,p,n,\text{pattern})$ 来自配置候选集，\texttt{seed} 用于可复现。

\paragraph{策略模式}
\begin{itemize}
\item \texttt{random}：随机采样动作（对照基线）。
\item \texttt{ucb}：基于历史 \texttt{objective} 的 UCB 采样（bandit 选择器）。
\item \texttt{pg}：softmax 参数化 policy-gradient bandit（带 baseline 的策略更新，状态持久化到 \texttt{outer\_policy\_pg.json}）。
\end{itemize}

\paragraph{内层执行模式}
当前外层每轮调用内层 \texttt{train\_only}，已支持两种停止方式：
\begin{itemize}
\item \texttt{fixed\_n}：按预算表数停止（推荐用于快速 outer-inner 交互）。
\item \texttt{converge}：沿用内层既有收敛逻辑（用于长训对照）。
\end{itemize}
其中 \texttt{fixed\_n} 已验证支持 checkpoint 正常落盘，并记录 \texttt{processed\_tables/stop\_reason/path\_read\_rate}。

\paragraph{外层问题形态与建模选择}
结合当前工程实现，外层每轮执行 ``一次动作 $\rightarrow$ 一次内层短训 $\rightarrow$ 一个标量回报''，且内层参数会持续变化，环境呈现\textbf{非平稳}。因此更适合建模为\textbf{非平稳 contextual bandit}，而非长时序 MDP。
这也是当前推荐优先采用 \texttt{ucb/pg}（以及后续可扩展的 TS）而不是重型 PPO/DQN 的核心原因：样本效率更高、调参更稳、与离散动作键天然匹配。

\paragraph{上下文状态定义（不新增数据源）}
外层状态建议仅由现有日志字段构成，避免引入新采集链路。定义：
\begin{equation}
s_t=\big[O_{t-1},J_{t-1},\rho_{t-1},q_{t-1},u_{t-1},\Delta O_{t-1},\kappa_t\big],
\end{equation}
其中 $O=\texttt{objective\_score}$，$J=\texttt{avg\_reward}$，$\rho=\texttt{action1\_rate}$，
$q=\texttt{path\_read\_rate}$，$u=\texttt{processed\_tables}/\texttt{table\_budget\_n}$，
$\Delta O_t=O_t-O_{t-1}$，$\kappa_t=t/T$（迭代进度）。实现时对各分量做滑窗标准化即可。

\paragraph{动作编码与规模控制}
当前动作为
\begin{equation}
a_g=(\mu_A,\mu_B,p,n,\text{pattern},\text{seed}),
\end{equation}
离散动作总规模为
\begin{equation}
|\mathcal{A}|=|\mathcal{M}_A|\cdot|\mathcal{M}_B|\cdot|\mathcal{P}|\cdot|\mathcal{N}|\cdot|\mathcal{T}|.
\end{equation}
其中 \texttt{seed} 建议作为\textbf{复现实验控制量}而非学习维度（避免把随机性误当作策略增益）。

\paragraph{奖励稳健化（建议）}
使用现有 \texttt{objective\_score} 作为外层即时奖励 $r_t$，并进行稳健处理：
\begin{enumerate}
\item 对有效轮的 $r_t$ 做滑窗标准化或截断（如 clip 到 $[-2,2]$）以抑制极端噪声；
\item 未通过 valid iter 门禁时，默认\textbf{不更新策略}（仅记录日志）；
\item 若 \texttt{path\_read\_rate} 低于阈值，可附加固定惩罚，保证策略不会朝坏数据收敛。
\end{enumerate}

\paragraph{策略实现细化与推荐}
\begin{itemize}
\item \texttt{random}（已实现）：仅作下界对照，不参与收敛判断的主结论。
\item \texttt{ucb}（已实现，推荐首选）：建议采用折扣统计应对非平稳，
\begin{equation}
Q_t(a)=\lambda_Q Q_{t-1}(a)+(1-\lambda_Q)\hat r_t(a),\quad
N_t(a)=\lambda_N N_{t-1}(a)+\mathbf{1}[a_t=a],
\end{equation}
\begin{equation}
\mathrm{score}_t(a)=Q_t(a)+c\sqrt{\frac{\log(\sum_{a'}N_t(a')+1)}{N_t(a)+\epsilon}}.
\end{equation}
\item \texttt{pg}（已实现）：softmax bandit 建议使用温度 $\tau$、EMA baseline 与熵正则，
\begin{equation}
\pi_\psi(a|s)=\mathrm{softmax}(z_a/\tau),\quad
\psi\leftarrow\psi+\alpha(r_t-b_t)\nabla_\psi\log\pi_\psi(a_t|s_t)+\beta\nabla_\psi\mathcal{H}(\pi_\psi).
\end{equation}
\item \texttt{ts}（建议后续增补）：可引入折扣 Thompson Sampling 作为高样本效率主力策略，优先级高于引入重型 actor-critic。
\end{itemize}

\paragraph{Phase2 / Phase3 的策略节奏建议}
\begin{itemize}
\item Phase2（外层收敛阶段）：保持较强探索（较大 $c$ 或较高 $\tau$），每个有效轮都更新外层策略。
\item Phase3（内层收敛阶段）：外层改为慢更新或间隔更新（每 $K$ 轮更新一次），防止外层频繁抖动干扰内层收敛。
\item 当 Phase3 进入后半段（接近 max=200）且指标稳定，可临时冻结外层策略，只验证内层是否达到切换条件。
\end{itemize}

\paragraph{默认超参数建议（起步值）}
\begin{itemize}
\item \texttt{ucb}: $\lambda_Q=\lambda_N=0.98,\ c=1.0,\ \epsilon=10^{-6}$。
\item \texttt{pg}: $\alpha=0.01,\ \beta=0.01,\ \tau=1.0$（可线性退火到 0.6），baseline EMA 系数 0.9。
\item valid iter 门禁保持本文既有规则，不额外新增判定字段。
\end{itemize}

\subsection{阶段切换规则（仅基于现有日志字段）}
\paragraph{原则}
不新增字段，不引入额外标注；仅使用现有 \texttt{outer\_train\_round.csv} 中的
\texttt{objective\_score, avg\_reward, action1\_rate, path\_read\_rate, processed\_tables, table\_budget\_n, stop\_reason, policy\_entropy} 进行判定。

\paragraph{通用有效轮（valid iter）门禁}
仅当满足以下条件时，该轮计入收敛统计：
\begin{enumerate}
\item \texttt{stop\_reason=fixed\_budget\_n}
\item \texttt{processed\_tables == table\_budget\_n}
\item \texttt{path\_read\_rate >= 0.99}
\end{enumerate}
未通过门禁的轮次仅记录，不参与阶段切换判据。

\paragraph{Phase2 $\rightarrow$ Phase3（外层 RL 收敛切换）}
设第 $t$ 轮有效轮的指标为 $O_t=\texttt{objective\_score}$，收敛条件定义为：
\begin{equation}
\left|\overline{O}_{t-4:t}-\overline{O}_{t-9:t-5}\right|\le \epsilon_O,\quad
\mathrm{Std}(O_{t-4:t})\le \sigma_O,
\end{equation}
建议默认 $\epsilon_O=0.01,\ \sigma_O=0.03$；若为 \texttt{pg} 模式再增加
\begin{equation}
\overline{\texttt{policy\_entropy}}_{t-4:t}\ge h_{\min},
\end{equation}
建议 $h_{\min}=0.05$，用于防止外层过早塌缩。\\
当上述条件\textbf{连续 3 次}满足时，判定 Phase2 收敛并切换到 Phase3。

\paragraph{Phase3 $\rightarrow$ Phase4（内层 RL 在外层数据上收敛切换）}
按你的约束固定为：\textbf{min=30, max=200}（以有效轮计数）。\\
设 $J_t=\texttt{avg\_reward}$，$\rho_t=\texttt{action1\_rate}$，收敛条件为：
\begin{equation}
\left|\overline{J}_{t-4:t}-\overline{J}_{t-9:t-5}\right|\le \epsilon_J,\quad
\mathrm{Std}(J_{t-4:t})\le \sigma_J,\quad
\overline{\rho}_{t-4:t}\ge \rho_{\min},
\end{equation}
建议默认 $\epsilon_J=0.01,\ \sigma_J=0.02,\ \rho_{\min}=0.02$。\\
当有效轮数 $\ge 30$ 且上述条件\textbf{连续 3 次}满足时，切换到 Phase4；若有效轮达到 200 仍未满足，则强制切换并标记 ``max\_reached\_not\_converged''。

\subsection{下一步执行清单（进入 S4 前）}
\begin{enumerate}
\item 完成 \texttt{random/ucb/pg} 的同口径对比跑（统一 action space、统一 seed 组、统一 \texttt{fixed\_n}）。
\item 固化 S4 课程超参（$\alpha$ 调度、replay 容量、采样比）并定义回退策略。
\item 按本节 ``阶段切换规则'' 固化自动切换阈值，并在 run 结束摘要中输出 ``是否达成连续3次'' 的判定结果。
\item 预生成 S5 实验清单与命名规范，避免后续结果文件不可比。
\end{enumerate}

\subsection{命令行执行蓝图（Windows，含可视化核查点）}
以下为\textbf{规划层面的命令模板}，用于后续实现时的统一运行口径。核心思想是：每条命令都能在控制台打印出 ``当前轮、当前路径、当前映射、当前训练摘要''，便于人工快速判断是否跑偏。

\paragraph{步骤 1：设置运行环境变量（一次性）}
\begin{verbatim}
# PowerShell
$env:RUN_ID="v3_edrl_o10_90_s42"
$env:RUN_ROOT="A:\MYpython\34959_RL\codes\logs\$env:RUN_ID"
$env:DYNAMIC_DATA_ROOT="$env:RUN_ROOT\post_stage\outer_batches\iter_001\R30"
$env:RL_DYNAMIC_INDEX_MODE="mod"
$env:RL_DYNAMIC_FILE_COUNT="12"
$env:RL_DYNAMIC_TABLE_BASE="400"
$env:RL_DYNAMIC_STRICT_PATH="1"
\end{verbatim}
预期控制台关键输出：
\begin{verbatim}
[OUTER][INIT] run_id=v3_edrl_o10_90_s42 ...
[OUTER][INIT] index_mode=mod file_count=12 table_base=400 strict=1
\end{verbatim}

\paragraph{步骤 2：执行单轮小批量生成（外层动作落地为文件）}
\begin{verbatim}
python codes\generation\generate_outer_small_batch.py `
  --run-id $env:RUN_ID `
  --iter-id 1 `
  --request-number 30 `
  --mu-a 10 --mu-b 90 --ratio-a 0.30 `
  --num-files 12 --seed 42001 `
  --out-root "$env:RUN_ROOT\post_stage\outer_batches"
\end{verbatim}
预期控制台关键输出：
\begin{verbatim}
[OUTER][GEN] iter=001 start ...
[OUTER][GEN] iter=001 done files=12 ok=12 fail=0
[OUTER][GEN] iter=001 stats ratio_A~0.30 ratio_B~0.70
\end{verbatim}
并且应生成：
\begin{verbatim}
post_stage/outer_batches/iter_001/R30/*.xlsx
post_stage/outer_generation.csv
post_stage/manifest.json
\end{verbatim}

\paragraph{步骤 3：执行内层短训（train\_only）并记录路径映射}
\begin{verbatim}
python codes\Dynamic_master34959.py `
  --dist_name O_10_90 `
  --request_number 30 `
  --algorithm PPO_NEW `
  --algo_version v3 `
  --stage-mode train_only `
  --seed 42 `
  --run-name $env:RUN_ID
\end{verbatim}
预期控制台关键输出：
\begin{verbatim}
[OUTER][PATH] iter=001 sample_map: t=400->0 ... t=411->11 t=412->0
[OUTER][PATH] iter=001 strict_check missing=0 unreadable=0
[OUTER][TRAIN] iter=001 learn_steps=... avg_reward=... action1_rate=...
\end{verbatim}

\paragraph{步骤 4：执行日志验收（决策对齐 + 表征对齐）}
\begin{verbatim}
python codes\analysis\validate_decision_artifacts.py `
  --run-dir "$env:RUN_ROOT" `
  --key-groups A B `
  --min-pairing-rate 0.98 `
  --min-h-align-rate 0.90 `
  --min-group-samples 30
\end{verbatim}
预期输出包含：
\begin{verbatim}
pairing_rate_matched >= 0.98
h_align_rate >= 0.90
[PASS] decision artifacts validated.
\end{verbatim}

\paragraph{步骤 5：多轮迭代时的最小巡检口径}
每轮迭代至少人工核查以下四行是否出现且数值合理：
\begin{itemize}
\item \texttt{[OUTER][ITER]}：动作参数是否变化（防止外层策略退化为常量）。
\item \texttt{[OUTER][GEN]}：文件数是否等于目标 \texttt{num\_files}，失败是否为 0。
\item \texttt{[OUTER][PATH]}：映射是否覆盖 ``跨界回绕''（如 11$\to$0）。
\item \texttt{[OUTER][TRAIN]}：\texttt{avg\_reward/action1\_rate} 是否可计算且非 NaN。
\end{itemize}

\subsection{最小验收阈值（建议）}
\begin{itemize}
\item 生成文件完整率：$=100\%$（严格模式下任何缺失直接失败）。
\item 路径映射读通率：$\ge 99\%$（建议目标 $100\%$）。
\item 决策日志配对率（pairing）：$\ge 98\%$。
\item h 对齐率（h\_align）：$\ge 90\%$。
\item 每轮 outer 产物完整性：\texttt{manifest + actions + generation + train\_round + path\_map} 五件套齐全。
\end{itemize}

\section{建议的对比实验与消融}
为使论文更有说服力，建议至少包含：
\begin{itemize}
\item 对比：V3（基线） vs V3+EDRL（完整）
\item 消融 A：去掉 $\Delta J$（纯对抗）\;vs\; 完整（验证``可学习性''必要性）
\item 消融 B：去掉反塌缩项 $B(E)$（验证动作塌缩机理）
\item 消融 C：Phase3 不混合 base 分布（验证课程混合/稳定性）
\item 诊断：solve-all/solve-none 类似指标（训练中 ``无信息'' 情形占比）与 action1\_rate 轨迹
\end{itemize}

\section{结论与下一步}
V3+EDRL 的价值在于：在不大改 V3 模型结构的情况下，把 OOD 攻坚点从``更强网络''转为``更有效的训练分布与样本结构''。外层策略通过事件文件参数化，完全兼容你们已有工程体系；通过 ``难但可学'' 奖励与反塌缩约束，避免对抗退化与样本饥饿。下一步最关键的工程验证是：外层是否能稳定学到一组 $a_g$ 使 O\_* 的 avg\_reward 与 action1\_rate 同时抬升，并且不显著牺牲 base 分布与 F\_*。

\section{NOVA\_EDRL 一体化主链路（2026-03-02 实现口径）}
\subsection{统一入口与分流规则}
\begin{itemize}
\item 主入口保持在 \texttt{codes/Dynamic\_master34959.py}。
\item 当 \texttt{--algorithm NOVA\_EDRL} 时，强制执行四阶段编排，不走旧单阶段分支。
\item 当算法为 \texttt{PPO\_NEW/PPO/A2C/...} 时，维持原有逻辑不变。
\end{itemize}

\subsection{四阶段执行语义（已落地）}
\begin{enumerate}
\item \textbf{Phase1（内层预训练）}：调用 master \texttt{train\_only}，生成/读取 1000 文件数据根，保存 \texttt{theta\_phase1}。
\item \textbf{Phase2（外层收敛）}：调用 \texttt{codes/outer\_rl/run\_edrl\_phase2.py}，外层策略在动作空间
\[
a_g=(\mu_A,\mu_B,p,n,\text{pattern},\text{seed})
\]
上选择，驱动 ``生成$\rightarrow$内层短训$\rightarrow$outer reward'' 闭环。
\item \textbf{Phase3（联合课程）}：继续使用同一外层主循环，启用课程混合与 replay；切换规则为有效轮口径下的收敛窗，且遵循 \texttt{min=30,max=200} 与连续满足约束。
\item \textbf{Phase4（复用 master implement）}：\textbf{不新建评测器}，直接回到 master，使用
\texttt{--stage-mode eval\_only --skip-generator --external-data-root <phase1\_data\_root> --init-model-path <phase3\_last\_ckpt>}，
复用既有 implement/test 评测链路与产物。
\end{enumerate}

\subsection{当前外层策略实现状态}
\begin{itemize}
\item 已实现三种可运行策略：\texttt{random}、\texttt{ucb}、\texttt{pg}。
\item \texttt{random} 用作下界与连通性烟测；\texttt{ucb}/\texttt{pg} 用于外层可学习策略对照。
\item 外层目标函数按统一日志字段计算（\(J,dJ,B,D\)），并支持路径严格校验与失败显式报错。
\end{itemize}

\subsection{实验入口（统一批跑）}
\begin{itemize}
\item \texttt{codes/experiments/run\_experiments\_server\_unified.py} 已支持
\texttt{--variant NOVA\_EDRL:v1}。
\item \texttt{run\_experiments\_common.py} 已将 \texttt{NOVA\_EDRL} 版本参数透传到 master，并调整运行键提取策略（优先 run name），避免内层 \texttt{meta.json} 覆盖导致的 NOVA 任务识别偏差。
\end{itemize}

\subsection{Phase4 产物口径}
Phase4 与旧 master implement 完全同口径，沿用现有评测输出（如 \texttt{rl\_trace.csv}、\texttt{rl\_training.csv}、\texttt{rl\_summary.csv} 及相关派生报表），因此不存在新的 phase4 专用依赖脚本。

\section{NOVA\_EDRL 最新流程与目标函数口径（2026-03-03）}
\subsection{四阶段完整流程（实现口径）}
\begin{enumerate}
\item \textbf{Phase1（内层预训练）}：在 base 分布上执行 master \texttt{train\_only}，产出并固化
\texttt{theta\_phase1.zip}、训练数据根目录与基础日志。
\item \textbf{Phase2（外层学习“出难题”）}：外层策略选择动作
\(
a_t=(\mu_a,\mu_b,p,n,\text{pattern},\text{seed})
\)
生成小批量事件文件；内层每轮都从同一个 \(\theta_{\mathrm{phase1}}\) 启动（冻结，不跨轮继承），据此评估外层奖励并更新外层策略参数。
\item \textbf{Phase3（联合训练）}：继续外层动作驱动数据生成，但内层恢复跨轮继承更新（第 \(t\) 轮的 \texttt{ckpt\_out} 作为第 \(t+1\) 轮 \texttt{ckpt\_in}），并执行课程混合与 replay。
\item \textbf{Phase4（复用 master implement）}：不新建评测器，直接回到 master 的 \texttt{eval\_only} implement 路径完成最终评测。
\end{enumerate}

\subsection{Phase2 目标函数（冻结内层）}
Phase2 的核心目标是“让外层学会打中内层弱点”，而不是让内层在该阶段持续变强。设
\(
J_t
\)
为第 \(t\) 轮在冻结内层参数下的平均回报（越大越容易），则难度项为
\(
1-J_t
\)。

最新 Phase2 外层目标定义为：
\begin{equation}
\label{eq:phase2_obj_latest}
\mathcal{O}^{(2)}_t
=
w_h\,(1-J_t)
+w_d\,\max(0, J_{\mathrm{ref}}-J_t)
+w_p\,H(a_t)
+w_m\,R_{\mathrm{minor}}(t)
-w_{\mathrm{path}}\,D_t
-w_{\mathrm{hard}}\,\max(0,J_{\mathrm{low}}-J_t)^2.
\end{equation}
其中：
\begin{itemize}
\item \(H(a_t)\)：由 \((\mu_a,\mu_b,p)\) 归一化得到的拥堵强度代理；
\item \(D_t\)：路径缺失/读取失败惩罚；
\item \(J_{\mathrm{ref}}\)：Phase2 参考回报（可固定或首轮自动初始化）；
\item \(J_{\mathrm{low}}\)：过难阈值（防止生成不可学样本）。
\end{itemize}

\paragraph{新增：少数动作恢复奖励}
先统计 Phase1 的两种动作占比，记少数动作为 \(a_{\min}\in\{0,1\}\)，其基线占比为
\(
\rho^{\mathrm{phase1}}_{\min}
\)。
在第 \(t\) 轮 Phase2 中，计算当前训练结果里该同一动作的占比
\(
\rho^{(t)}_{\min}
\)，定义：
\begin{equation}
R_{\mathrm{minor}}(t)=\max\!\left(0,\rho^{(t)}_{\min}-\rho^{\mathrm{phase1}}_{\min}\right).
\end{equation}
即：若 Phase1 少数动作占比被拉升，则给予外层正奖励；未拉升则不加分。

\subsection{Phase3 目标函数（保持不变）}
Phase3 继续使用既有 EDRL 目标，不做改动：
\begin{equation}
\mathcal{O}^{(3)}_t
=
J_t+\lambda_{\Delta J}\,dJ_t
B_t
-\gamma_{\mathrm{path}}D_t
-\text{floor\_soft}_t
-\text{floor\_hard}_t.
\end{equation}
其中 \(B_t\) 与 floor 项负责反塌缩约束，\(dJ_t\) 反映跨轮改进。

\subsection{参数继承与切换规则（最新）}
\begin{itemize}
\item \textbf{Phase2}：\texttt{ckpt\_in} 固定为 \(\theta_{\mathrm{phase1}}\)，\texttt{ckpt\_out} 仅用于审计，不用于下一轮 Phase2 初始化。
\item \textbf{Phase3}：恢复跨轮继承，\texttt{ckpt\_in}=\texttt{prev ckpt\_out}。
\item \textbf{Phase2\(\rightarrow\)Phase3}：沿用外层收敛窗规则（有效轮口径，默认 \texttt{min=10,max=80}）。
\item \textbf{Phase3\(\rightarrow\)Phase4}：有效轮 \(\ge 30\) 且收敛条件连续满足 3 次即切换；若达到 \texttt{max=200} 仍未满足则强制切换。
\end{itemize}

\subsection{日志可验证性（关键字段）}
\begin{itemize}
\item \texttt{outer\_train\_round.csv}: \texttt{phase, ckpt\_in, ckpt\_out, objective\_score, avg\_reward, action0\_rate, action1\_rate, minority\_rate, stop\_reason, processed\_tables, path\_read\_rate}
\item Phase2 诊断打印：\texttt{[OUTER][P2-OBJ] hard/drop/proxy/minority\_gain/minority\_reward/too\_hard}
\item \texttt{post\_stage/phase2\_action\_baseline.json}: 保存 Phase1 动作占比基线与少数动作标识，供 Phase2 直接复用。
\end{itemize}

\subsection{2026-03-03 Implementation Delta (CQL/CaDM)}
\begin{itemize}
\item \textbf{CQL diagnostics completed}: \texttt{codes/robust\_rl/cql\_dqn.py} now exports
\texttt{cql\_alpha,cql\_temp,cql\_updates,cql\_td\_loss,cql\_cql\_loss,cql\_q\_mean,cql\_q\_std,cql\_q\_max,cql\_q\_taken,cql\_lse\_q,cql\_ood\_q\_gap}.
The fields are persisted through \texttt{codes/core/dynamic\_RL34959.py} into \texttt{rl\_training.csv}.
\item \textbf{CQL alpha ablation script hardened}: \texttt{codes/analysis/run\_cql\_alpha\_ablation.py}
now selects the last \emph{metric-bearing} train row (instead of the last train row blindly), and falls back to the last eval row for
\texttt{avg\_reward}/\texttt{rolling\_avg} when needed.
\item \textbf{CaDM master routing fixed}: \texttt{codes/Dynamic\_master34959.py} now routes
\texttt{--algorithm CADM} to \texttt{PPO\_NEW} and defaults to \texttt{algo\_version=v6.3\_cadm}
when the caller does not explicitly override version.
\item \textbf{Smoke verification status}: \texttt{smoke\_cadm\_master\_default} confirms
\texttt{cadm\_enabled=1} and non-empty auxiliary losses in \texttt{rl\_training.csv}.
\item \textbf{RARL adversary refinement}: \texttt{run\_edrl\_phase2.py} now supports
strict zero-sum reward for adversary updates (\texttt{--rarl-zero-sum-strict=1} uses \(-J\)),
minimum replay gate (\texttt{--rarl-min-replay}), and windowed adversary state features
(recent \(J\), action ratio, entropy via \texttt{--rarl-state-window}).
Policy-mode/objective-mode mismatch can be auto-corrected with \texttt{--rarl-force-objective=1}.
\item \textbf{PLR plot-ready summary added}: \texttt{codes/analysis/summarize\_outer\_plr\_stats.py}
exports per-run curve CSV (\texttt{outer\_plr\_curve.csv}), top-priority level table
(\texttt{outer\_plr\_top\_levels.csv}), and multi-run summary
(\texttt{plr\_replay\_experiment\_summary.csv}) with replay/new ratios, top-k coverage, and
priority-sample share for direct paper plotting.
\item \textbf{Master baseline wiring corrected}: \texttt{codes/Dynamic\_master34959.py}
now maps \texttt{RARL} profile to \texttt{--outer-policy-mode rarl\_dqn --objective-mode rarl}
(with strict zero-sum defaults), and maps \texttt{PLR\_UED} profile to true level-replay mode
(\texttt{--plr-level-replay} plus \texttt{p\_new}/buffer/EMA/min-weight defaults), so the two
reference baselines execute their core mechanisms by default.
\end{itemize}

\appendix
\section{外层动作空间的两个实用版本}
\subsection{版本 A：离散 key（推荐先做）}
外层动作直接从你们现有 distribution\_config 的 key 集采样，$a_g\equiv k$。优点：实现最省事、可复现、审稿友好；缺点：探索粒度受限。

\subsection{版本 B：参数版（更灵活）}
按式\eqref{eq:ag} 输出参数，再映射为事件文件。建议先只放开 $(\mu_1,\mu_2,p,\text{pattern})$，其余固定。

\section{可直接复用的记录字段（meta.json）}
建议每次 run 记录：
\begin{itemize}
\item algo\_version: v3+edrl
\item phase: 1/2/3/4
\item generator\_action: $(\mu_1,\mu_2,p,\text{pattern},\sigma,\text{seed})$
\item events\_file\_path / hash
\item ppo\_hyperparams / generator\_hyperparams
\item summary: avg\_reward, action1\_rate, p\_action1, entropy, reward\_var
\end{itemize}

\end{document}
